\section{Introduction}

The generalized minimal residual method (GMRES) constructs approximations in a
Krylov subspace and minimizes a residual norm through a small least-squares
problem \cite{saad_schultz_1986}. Its use for nonsymmetric and complex-valued
systems makes implementation details---orthogonalization, least-squares updates,
stopping rules, and restarting---numerically consequential. The Arnoldi process
underlying GMRES originates in the construction of orthogonal Krylov bases
\cite{arnoldi_1951}; standard algorithmic treatments are given by
Saad \cite{saad_2003}.

This article records the first numerical observations from the WaveCore Research
GMRES implementation. The scope is deliberately narrow: deterministic dense test
matrices are converted to the project's complex CSR format and solved by the same
public routine intended for later sparse applications. Internal instrumentation
observes the actual Arnoldi basis without changing the returned solution or
status. The objectives are to verify implementation-level invariants, identify
failure modes before coupling the solver to a Helmholtz discretization, and retain
a traceable account of measured values.

No new GMRES variant or general stability theorem is proposed. In particular,
successful behavior on these small systems does not establish robustness for all
matrix classes, dimensions, restart lengths, or floating-point environments.
